\chapter{Statistics}\label{ch:iv:statistics}

``Explain a p-value to a trader.'' The candidate says: ``It is the probability that the strategy does not work.''
The interviewer writes one word in the margin and asks the next question, which is the same question about a
Sharpe ratio of 2 measured over nine months. Statistics questions in interviews look elementary and are not:
they test whether the candidate knows what an estimate's error is, how a regression can mislead, what a test does
and does not say, and how to turn a claim made on a desk into something that can be checked. The mathematics is
One Quant Book~4's (chapters 11 to~16); this chapter is about answering with it.

\section{Estimation and the error of what you report}

Every number reported from data has a standard error, and a good answer says it. Two cases come up constantly.
The mean return: with daily volatility $\sigma$ and $n$ days, the mean daily return has standard error
$\sigma/\sqrt n$, so a year of data estimates the annual mean return with a standard error equal to the annual
volatility itself. And the Sharpe ratio (One Quant Book~4, chapter~11): for independent returns its standard
error is approximately $\sqrt{(1 + \mathrm{SR}^2/2)/n}$ in per-period units (Lo, 2002).

\begin{example}[A Sharpe ratio of 2 over nine months]\label{ex:iv:statistics:sharpe}
Nine months are about 189 trading days. A daily Sharpe ratio of $2/\sqrt{252} \approx 0.126$ has a standard error
of $\sqrt{(1 + 0.126^2/2)/189} \approx 0.073$, which annualises to about 1.16. The estimate is $2.0 \pm 1.16$: a
$t$-statistic of 1.73, not significant at 5\% on a two-sided test, before any adjustment for how many strategies
were tried or for autocorrelation. A desk that allocates on nine months of Sharpe ratio is allocating mostly on noise.
\end{example}

\begin{method}[Answering any ``is this estimate good?'' question]\label{met:iv:statistics:estimate}
\begin{enumerate}
\item Say what is being estimated, from how many effectively independent observations.
\item Give the standard error, from a formula or a bootstrap that respects the data's dependence (One Quant
Book~4, chapter~13).
\item Say what would bias the estimate: selection of the best of many, look-ahead, survivorship, overlapping
observations.
\item Give the interval and the decision it supports.
\end{enumerate}
\end{method}

\section{Regression and its pathologies}

Least squares (One Quant Book~4, chapter~16) is the tool of most research answers, and interviews probe its failure
modes.

\begin{itemize}
\item \emph{Two regressions.} The slope of $y$ on $x$ times the slope of $x$ on $y$ is $r^2$, so the two lines
coincide only when $|r| = 1$; ``regression to the mean'' is this asymmetry.
\item \emph{Omitted variables.} If $y = \beta_1 x_1 + \beta_2 x_2 + \varepsilon$ and $x_2$ is left out, the slope
on $x_1$ estimates $\beta_1 + \beta_2 \Cov(x_1, x_2)/\Var(x_1)$, which can have the opposite sign to $\beta_1$.
\item \emph{Errors in variables.} Noise in the regressor attenuates the slope by $\Var(x)/(\Var(x) + \Var(u))$; a
signal measured with as much noise as signal has its coefficient halved (attenuation bias).
\item \emph{Collinearity.} Near-collinear regressors leave the fit good and the individual coefficients unstable;
ridge regression trades bias for variance.
\item \emph{Outliers.} One extreme day can make a regression; always ask what the result is without the largest
five observations.
\end{itemize}

\begin{example}[A sign that flips]\label{ex:iv:statistics:omitted}
Returns load $+1$ on a value signal $x_1$ and $+2$ on a momentum signal $x_2$, and the two signals have correlation
$-0.8$ with unit variances. Regressing returns on value alone gives a slope of $1 + 2 \times (-0.8) = -0.6$: value
appears to lose money because it is short momentum. A simulation of 200\,000 observations gives the same.
\end{example}

\section{Testing}

A p-value is the probability, computed under the null hypothesis, of a test statistic at least as extreme as the one
observed (One Quant Book~4, chapter~12). It is not the probability that the null is true, not the probability that the
result is due to chance, and not a measure of the effect's size; the American Statistical Association's statement
(Wasserstein and Lazar, 2016) lists these misreadings. Three further facts carry most interview questions.

\begin{itemize}
\item \emph{Power.} To detect an annual Sharpe ratio of $s$ with a one-sided 5\% test and 80\% power needs about
$\big((1.645 + 0.842)/s\big)^2$ years: 25 years for $s = 0.5$, six for $s = 1$.
\item \emph{Many tests.} Among $m$ true nulls tested at level $\alpha$, $m\alpha$ are expected to be rejected; the best
of $m$ backtests is biased upward (\cref{prop:iv:applications:max}), and the deflated Sharpe ratio corrects for it.
\item \emph{Stopping rules.} The p-value depends on the experiment that would have been run had the data differed,
not only on the data; the same heads and tails give different p-values under ``toss twelve times'' and ``toss until
three heads'' (Lindley and Phillips, 1976, make the point with their own numbers).
\end{itemize}

\begin{center}
\small
\begin{tabular}{@{}lccccc@{}}
\toprule
Backtests tried, all worthless & 1 & 10 & 100 & 200 & 1\,000 \\
\midrule
Expected best annual Sharpe ratio over 5 years & 0 & 0.69 & 1.12 & 1.23 & 1.45 \\
\bottomrule
\end{tabular}

\smallskip
{\small The expected best of $k$ independent zero-skill Sharpe ratios estimated over five years:
$\E[\max_{i \le k} Z_i]/\sqrt5$. Correlated trials give a smaller but still positive bias.}
\end{center}

\Cref{fig:iv:statistics:power} turns the power calculation of \cref{iq:iv:statistics:8} into a chart: the years of daily
returns needed before a one-sided 5\% test detects a given annual Sharpe ratio.

\begin{omfigure}
\begin{tikzpicture}
\begin{semilogyaxis}[width=9.5cm,height=5.4cm,xlabel={\small true annual Sharpe ratio},ylabel={\small years of data},
  xmin=0.25,xmax=3,ymin=0.1,ymax=200,
  tick label style={font=\footnotesize},grid=major,grid style={gray!20},table/col sep=comma,
  legend cell align=left,legend pos=north east,legend style={font=\scriptsize}]
\addplot[omDef,thick] table[x=sr,y=p50]{figdata/interviews/14-statistics/power.csv};
\addplot[omProp,thick,dashed] table[x=sr,y=p80]{figdata/interviews/14-statistics/power.csv};
\addplot[omThm,thick,densely dotted] table[x=sr,y=p95]{figdata/interviews/14-statistics/power.csv};
\legend{50\% power,80\% power,95\% power}
\end{semilogyaxis}
\end{tikzpicture}

\omcaption{Years of returns needed for a one-sided 5\% test to detect a strategy with a given true annual Sharpe ratio, at three
levels of power, from $T = ((z_{0.95} + z_{\text{power}})/\mathrm{SR})^2$. A Sharpe ratio of 1 needs about six years for 80\%
power, one of 0.5 about 25. Data: \texttt{fig\_iv\_power.py}.}\label{fig:iv:statistics:power}
\end{omfigure}

\section{Turning a desk claim into a test}

Interviewers often state a claim in desk language and wait for the candidate to formalise it.

\begin{method}[From a claim to a test]\label{met:iv:statistics:claim}
\begin{enumerate}
\item Write the claim as a statement about a parameter: ``the signal predicts'' becomes ``the slope of next-day
return on today's signal is positive''.
\item State the null and the alternative, and the unit of independent observation (a day, a week, a stock-day
cluster).
\item Choose the statistic and its standard error, robust to the heteroskedasticity and dependence the data have.
\item Say how many related claims were examined, and adjust.
\item Say what result would change your mind, before looking.
\end{enumerate}
\end{method}

\section{Worked answers}

\begin{example}[``Our fills are better on Mondays'']\label{ex:iv:statistics:monday}
Over a year, the desk's average slippage was 1.2 basis points on 52 Mondays and 1.8 on the other 208 days, with a daily
standard deviation of 1.5. Following \cref{met:iv:statistics:claim}: the parameter is the difference in mean slippage;
the unit is a day; the statistic is $0.6/(1.5\sqrt{1/52 + 1/208}) \approx 0.6/0.23 \approx 2.6$, a two-sided p-value of
about 0.01. Then step four: someone looked at five weekdays and reported the best one. With a Bonferroni correction the
p-value becomes about $5 \times 0.0099 \approx 0.05$, borderline at best, and the honest summary is ``worth checking on
next year's data, with Monday named in advance''. What would change the conclusion: a mechanism (a Monday auction, a
weekend news effect) that predicts the sign before the data are seen.
\end{example}

\begin{example}[A hit rate from clustered trades]\label{ex:iv:statistics:hit}
\emph{``540 of 1\,000 trades made money. Is the strategy better than a coin?''} Treating the trades as independent, the
standard error is $\sqrt{0.54 \times 0.46/1000} \approx 1.58\%$ and the 95\% interval $54\% \pm 3.1\%$ excludes 50\%. But
the trades came in 100 days of 10, and trades on the same day share the day's market move; with an intra-day correlation
of outcomes of 0.1, the variance is multiplied by the design effect $1 + (10 - 1) \times 0.1 = 1.9$, the standard error by
$\sqrt{1.9} \approx 1.38$, and the interval becomes $54\% \pm 4.3\%$, which includes 50\%. The answer names the unit of
independence before it names a p-value.
\end{example}

\begin{example}[Last year's top decile]\label{ex:iv:statistics:decile}
\emph{``The top tenth of a hundred traders by last year's P\&L are given more capital. If yearly performance has a
correlation of 0.3 from one year to the next, what do you expect from them this year?''} In standard units, the top
decile of a normal averages $\varphi(1.2816)/0.1 \approx 1.75$ standard deviations; the best prediction for next year is
$0.3 \times 1.75 \approx 0.53$. They are still expected to be above average, at less than a third of last year's
distance from it: regression to the mean, not a loss of skill. The same arithmetic applies to the best backtest of a
search (\cref{iq:iv:statistics:10}).
\end{example}

\section{Question bank}

\begin{interviewq}[$\star$ \iqroles{researcher, trader} \iqfirm{any}]\label{iq:iv:statistics:1}
Define a p-value in one sentence. Then say what is wrong with ``the probability that the strategy does not work''.
\end{interviewq}

\begin{interviewq}[$\star$ \iqroles{researcher, risk} \iqfirm{systematic fund}]\label{iq:iv:statistics:2}
Daily returns have a standard deviation of 1\%. With 250 days of data, what is the standard error of the mean daily
return, and of the annual mean return? What does this say about estimating expected returns?
\end{interviewq}

\begin{interviewq}[$\star$ \iqroles{researcher, mle} \iqfirm{any}]\label{iq:iv:statistics:3}
Why does the sample variance divide by $n - 1$? When does it matter in practice?
\end{interviewq}

\begin{interviewq}[$\star$ \iqroles{researcher} \iqfirm{systematic fund}]\label{iq:iv:statistics:4}
The slope of $y$ on $x$ is 0.8 and the slope of $x$ on $y$ is 0.45. What is the correlation between them?
\end{interviewq}

\begin{interviewq}[$\star\star$ \iqroles{trader, researcher} \iqfirm{multi-manager fund}]\label{iq:iv:statistics:5}
A strategy shows an annualised Sharpe ratio of 2 over nine months of daily returns. How precisely is that Sharpe ratio
known, and is it significantly different from zero?
\end{interviewq}

\begin{interviewq}[$\star\star$ \iqroles{researcher} \iqfirm{systematic fund}]\label{iq:iv:statistics:6}
Returns load $+1$ on signal $x_1$ and $+2$ on signal $x_2$; the signals have unit variance and correlation $-0.8$.
What slope do you get regressing returns on $x_1$ alone? What would you conclude, wrongly?
\end{interviewq}

\begin{interviewq}[$\star\star$ \iqroles{researcher, mle} \iqfirm{systematic fund}]\label{iq:iv:statistics:7}
Your signal is measured with noise whose variance equals the signal's. By what factor is its regression coefficient
biased, and in which direction? How would you correct it?
\end{interviewq}

\begin{interviewq}[$\star\star$ \iqroles{researcher, trader} \iqfirm{multi-manager fund}]\label{iq:iv:statistics:8}
How many years of returns do you need to tell a strategy with an annual Sharpe ratio of 0.5 from zero, with a
one-sided 5\% test and 80\% power? And for a Sharpe ratio of 1?
\end{interviewq}

\begin{interviewq}[$\star\star$ \iqroles{researcher, risk} \iqfirm{any}]\label{iq:iv:statistics:9}
You test 100 signals, all of them useless, each at the 5\% level. How many do you expect to find significant? What is
the chance of at least one? What threshold would Bonferroni use?
\end{interviewq}

\begin{interviewq}[$\star\star\star$ \iqroles{researcher} \iqfirm{systematic fund}]\label{iq:iv:statistics:10}
A colleague tried 200 variants of a strategy over five years and reports the best, with an annualised Sharpe ratio of
1.9. What would the best of 200 worthless variants show on average? How do you report the result?
\end{interviewq}

\begin{interviewq}[$\star\star\star$ \iqroles{researcher, mle} \iqfirm{systematic fund}]\label{iq:iv:statistics:11}
Daily P\&L follows an AR(1) with coefficient 0.5. By what factor does the naive standard error of the mean understate
the true one? How would you bootstrap it?
\end{interviewq}

\begin{interviewq}[$\star\star\star$ \iqroles{researcher, trader} \iqfirm{any}]\label{iq:iv:statistics:12}
A coin showed 3 heads in 12 tosses. Test ``the coin is fair'' against ``it favours tails'' if the plan was to toss 12
times, and again if the plan was to toss until the third head. Why do the p-values differ, and does it trouble you?
\end{interviewq}

\begin{interviewq}[$\star\star\star$ \iqroles{researcher, risk} \iqfirm{bank}]\label{iq:iv:statistics:13}
A regression of a stock's daily return on a factor over a year has $R^2 = 0.21$. On inspection, one day of the 251 was a
crash in which both fell twelve standard deviations. What do you expect $R^2$ to be without that day, and what do you
report?
\end{interviewq}

\begin{omsources}
\item A.~W.~Lo, ``The statistics of Sharpe ratios'', \emph{Financial Analysts Journal} 58(4), 2002, 36--52.
\item R.~L.~Wasserstein and N.~A.~Lazar, ``The ASA statement on p-values: context, process, and purpose'',
\emph{The American Statistician} 70(2), 2016, 129--133.
\item D.~V.~Lindley and L.~D.~Phillips, ``Inference for a Bernoulli process (a Bayesian view)'', \emph{The American
Statistician} 30(3), 1976, 112--119.
\item One Quant Book~4, chapters 11--16 (estimation, testing, resampling, linear models).
\end{omsources}
